\section{Experimental System: Tasks and the Blackboard Game}
\label{sec:methods_game}

Our experiment separates task construction from collective decision-making.
We first create a team-allocation problem with a known answer and divide its
evidence among language-model agents. We then let the agents exchange that
evidence and revise their answers through a shared message board. A controller
can post on the same board to encourage a chosen answer. The resulting game
lets us measure how communication changes both the information available to
agents and their collective decision.

\subsection{Constructing the task}
\label{sec:task_generation}

\paragraph{A team-allocation problem.}
Our tasks adapt the Team Allocation setting of MuSR
\citep{sprague2024musr}, which tests multistep reasoning from natural-language
evidence. Three people must be assigned to two jobs: one person works alone on
the first job, and the remaining pair works together on the second. There are
therefore three possible answers. For example, suppose the jobs are to analyze
field measurements and organize a community workshop. The choices are:
\begin{center}
\begin{tabular}{ll}
\toprule
Analyze measurements & Organize the workshop \\
\midrule
Alice & Bruno and Chandra \\
Bruno & Alice and Chandra \\
Chandra & Alice and Bruno \\
\bottomrule
\end{tabular}
\end{center}
Choosing among them requires evidence about individual skills and cooperation
within each possible pair. We generate our own instances of this task family.

\paragraph{Establishing the answer before writing the evidence.}
We first construct a hidden numerical model with six skill values (one per
person and job) and three pairwise cooperation values. Each value lies in
$\{1,2,3\}$. If person $p$ takes the first job and people $u,v$ take the second,
our allocation score is
\begin{equation}
    S(p;u,v)=s_{p,1}+s_{u,2}+s_{v,2}+c_{u,v},
\end{equation}
where $s_{p,t}$ measures person $p$'s skill for job $t$ and $c_{u,v}$ measures
cooperation. We evaluate all three allocations and retain worlds with a unique
highest-scoring answer. Thus, the correct answer is determined by code before
any language-model interaction. Agents receive the scenario and evidence,
while the numerical model and scores remain hidden.

\paragraph{Turning the hidden model into evidence cards.}
Following MuSR's approach of constructing a reasoning problem before rendering
it in language \citep{sprague2024musr}, we create short accounts of behavior
that support facts about skills or cooperation. For example, an account of
Alice identifying and repairing a spreadsheet error provides evidence of her
analytical ability. We call each short account an \emph{evidence card}.
Every card is linked to a proposition that can be checked against the hidden
model, such as a person's skill level or a comparison between two people's
skills. An LLM writes the accounts; equality comparisons use fixed templates.
We check that the text supports its intended proposition and does not disclose
the hidden scores or recommend an answer. This preserves the distinction
between the exact answer, established by the numerical model, and the
natural-language reasoning needed to infer it from the cards.

\paragraph{Distributing incomplete information.}
We assign one private card to each of $N=24$ agents. A single card leaves
several answers plausible, whereas combining cards provides stronger evidence.
For example, one agent may know about Alice's analytical ability and another
about Bruno and Chandra's cooperation. We assess this partial-information
ambiguity by enumerating the numerical models consistent with each private
card and checking which allocation wins in each model. The assignment includes
complementary cards that favor the correct answer when considered together.
Agents may initially guess correctly, but they must exchange evidence to
access information beyond their own card.

We freeze the task and evidence before play. Each agent first votes using only
its private card. These initial votes are saved and reused across matched
experimental conditions. For each model call, we shuffle the A/B/C labels
assigned to the three allocations and translate the response back to the
chosen allocation, preventing a fixed answer letter from becoming a shared
social cue.

\subsection{Playing the blackboard game}
\label{sec:relational_game}

The \emph{blackboard} is a shared collection of short-lived messages. Agents
take turns reading a small random sample of posts, choosing an answer, and
optionally adding a post of their own. A message posted early can therefore
influence later decisions within the same round. For example, an agent may
post the card about Alice's analytical ability; another agent who reads it
can combine it with its own cooperation evidence and reconsider the allocation.
Access to the board is limited: agents see only the messages they sample.

\paragraph{Agent state and memory.}
Agent $i$ maintains a current answer $X_i$, an active evidence set $K_i$, and a
historical set $H_i$ containing every card it has received. Only active evidence
is available in its prompt. At each round boundary, each active card survives
independently with probability $\rho$; historical records are retained.
Reading a previously lost card makes it active again. The parameter $\rho$
therefore controls how long communicated evidence remains usable. Previous
private explanations are recorded for analysis but are not fed back to agents.

\paragraph{One round.}
Each round has a boundary stage, a controller posting stage, and $N$ sequential
agent updates. The controller first samples votes and decides whether to act.
The previous round's messages then expire, evidence persistence is applied,
and an active controller posts its messages. During each subsequent update:
\begin{enumerate}[leftmargin=*,nosep]
    \item One focal agent is sampled uniformly from the population. It reads
    up to $q$ live board messages, sampled uniformly without replacement,
    excluding its own posts.
    \item The agent receives the question, answer options, active evidence,
    previous vote, and sampled messages. One LLM call returns an updated vote,
    a private explanation, and an optional public message.
    \item The vote is updated immediately, attached evidence from the sampled
    reports is added to the agent's active and historical sets, and any new
    public message becomes available to later updates in the same round.
\end{enumerate}
Focal agents are sampled independently at each position, so a round need not
update every agent once. Messages live only for their creation round; an empty
board supplies no social messages. Episodes run for a fixed number of rounds,
including after consensus.

\paragraph{Messages and evidence transmission.}
An ordinary agent may post a \texttt{REPORT}, ask for evidence with a
\texttt{REQUEST}, or post nothing. A report contains public text and may attach
one evidence identifier available to its author. Requests carry no evidence.
Posts also carry the author's current allocation choice, and replies may refer
only to messages visible in that update. The prompt distinguishes a
participant's interpretation from the canonical evidence text attached to a
report. Only reading a report with a valid attachment changes the exact
evidence state. Posting alone does not transmit a card to anyone, and public
prose is not counted as verified evidence. This separates the transmission of
task evidence from persuasion through an agent's interpretation of that evidence.

\paragraph{Feedback through the same public board.}
Let $Z$ be the controller's target, which may be correct or incorrect, and let
\begin{equation}
    n_k=\sum_{i=1}^{N}\mathbb{I}[X_i(k)=Z],
    \qquad x_k=n_k/N
\end{equation}
be its support before round $k$. The controller samples $q_c$ distinct agents
and observes their votes, giving a target-support count $Y_k$. The binary
feedback policy in Section~\ref{sec:feedback_efficiency} chooses between
silence ($U_k=0$) and communication ($U_k=1$); lower sampled target support
makes communication more likely. This decision uses votes only.

The controller has a fixed pool of true evidence cards. To favor an incorrect
answer, it selectively discloses evidence that keeps that answer plausible
while withholding evidence that would weaken it. Its reports remain truthful
at the fact level.

On active rounds, an adaptive communication policy chooses to publish up to
$b$ evidence reports, make one request for evidence, or post one
\texttt{DIRECTIVE} directing attention to a comparison. The latter two modes
carry no evidence. An LLM selects the mode and
eligible report identifiers using sampled votes, previous public messages,
the allowed fact pool, and posting history. It has no access to agents'
private evidence or explanations. The selected reports reproduce the frozen
evidence text, while requests and directives use fixed templates. This
communication choice occurs only after the binary decision to act. A
report-only variant instead posts exactly $b$ distinct evidence reports on
each active round.

The controller appears under a stable participant identity. Its messages
enter the same board and have no sampling priority. Consequently, $b$ bounds
posting capacity, while actual exposure depends on which messages agents
read. We record posts, reads, evidence acquisition and reactivation, and vote
changes separately. Tracking support for both $Z$ and the correct answer,
alongside active and historical evidence, allows us to distinguish successful
steering from improved task performance.
